\section*{Chapter \ref{ch:b3:trade-policy} --- Trade Policy}

\begin{solution}{exo:b3:trade-policy:1}
In points, India: from 56.4~\% in 1990 to 4.6~\% in 2022, 51.8 points. In percent,
China: from 32.2~\% to 2.2~\%, a cut of 93~\% (India's is 92~\%). % figdata: tariffs_points
China's fell most in the 1990s, from 32.2~\% in 1992 to 14.7~\% in 2000, and then to
4.7~\% by 2010. % figdata: tariffs
\end{solution}

\begin{solution}{exo:b3:trade-policy:2}
Free trade: demand 40, supply 20, imports 20. With the tariff, $P = 45$: demand 35,
supply 25, imports 10. Revenue $5 \times 10 = 50$. Deadweight loss $\tfrac12 \times 5 \times 5 +
\tfrac12 \times 5 \times 5 = 25$ (production and consumption triangles).
\end{solution}

\begin{solution}{exo:b3:trade-policy:3}
$a = 12{,}000/20{,}000 = 0.6$. ERP $= (0.10 - 0.6 \times 0.05)/0.4 = 0.07/0.4 = 17.5~\%$, nearly
twice the nominal tariff. With 20~\% on inputs: $(0.10 - 0.12)/0.4 = -5~\%$: the industry is
taxed, not protected.
\end{solution}

\begin{solution}{exo:b3:trade-policy:4}
Compare the prices foreign exporters receive, before the tariff (at the border,
tariff excluded), for the goods targeted against comparable goods not targeted, or
against the same goods sold to other countries. Full absorption means these
prices fell by the full tariff, so that the duty-inclusive price paid by importers
did not change.
\end{solution}

\begin{solution}{exo:b3:trade-policy:5}
Free trade: $100 - 2P = 2P - 20$, $P = 30$, $M = 40$. Tariff: $100 - 2(P^* + 10) = 2P^* - 20$,
so $P^* = 25$, $P = 35$, $M = 30$. Terms-of-trade gain $(30 - 25) \times 30 = 150$; deadweight
loss $\tfrac12 \times (35 - 30) \times (40 - 30) = 25$; net $+125$. (Check: revenue $10 \times 30 = 300$,
minus the loss of home surplus $5 \times (40 + 30)/2 = 175$.)
\end{solution}

\begin{solution}{exo:b3:trade-policy:6}
$t = 1/4 = 25~\%$. As the elasticity grows, the \omterm{def:b3:trade-policy:optimal-tariff}{optimal tariff} falls toward zero, the
small-country case.
\end{solution}

\begin{solution}{exo:b3:trade-policy:7}
Free trade for ever gives 0 each year. Deviating gives 4 once and $-2$ for ever
after: $4 - 2\beta/(1 - \beta)$ in present value. Free trade is an equilibrium when $0 \ge 4 -
2\beta/(1 - \beta)$, that is $4(1 - \beta) \le 2\beta$, or $\beta \ge 2/3$. Punishing by tariffs is
itself a Nash equilibrium of the stage game, so the threat is credible.
\end{solution}

\begin{solution}{exo:b3:trade-policy:8}
The tested solution is \code{ex08_overhang.py}, in the chapter's \code{exercises/}
folder. With income per head available, 128 countries; the median overhang is 20.9
points, and 82~\% of the countries could double their applied tariff within their
bindings. The overhang falls with income: the correlation with log income per head
is $-0.41$, and the median overhang is 32.6 points in the poorer half against 12.1 in
the richer half. % output: ex08_overhang
\end{solution}

\begin{solution}{exo:b3:trade-policy:9}
No: the industry might have grown without protection, as did others in the same
country, or grown in spite of it. Refutation: the protected industry's
productivity growth no faster than that of comparable unprotected industries, or
than the same industry abroad; its costs still above world costs when protection
ends; or the gains smaller than the cost of protection to buyers in the meantime.
A clean test needs protection that varies for reasons unrelated to the industry's
prospects, as in the Napoleonic blockade.
\end{solution}

\begin{solution}{exo:b3:trade-policy:10}
Home welfare is surplus plus revenue $sM$. With $dP = ds + dP^*$: $dW = -M\,dP + M\,ds +
s\,dM = -M\,dP^* + s\,dM$. Market clearing gives $dM = X^{*\prime}dP^*$, so $dW = 0$ when $-M + sX^{*\prime} =
0$: $s = M/X^{*\prime} = X^*/X^{*\prime}$, and $s/P^* = X^*/(P^*X^{*\prime}) = 1/\varepsilon^*$.
\end{solution}

\begin{solution}{exo:b3:trade-policy:11}
Tariff: imports available at 50; the firm produces where $20 + q = 50$, $q = 30$;
demand $100 - 50 = 50$; imports 20. Quota of 20: residual demand $P = 80 - q$, marginal
revenue $80 - 2q = 20 + q$, so $q = 20$ and $P = 60$; demand 40, imports 20. Same imports,
but the quota gives a price of 60 instead of 50, and domestic output of 20 instead
of 30: the quota creates market power, the tariff does not.
\end{solution}

\begin{solution}{exo:b3:trade-policy:12}
The tested solution is \code{ex12_shift_share.py}, in the chapter's
\code{exercises/} folder. Exposures: A 2.75, B 1.60, C 0.75, D 2.30, E 0.30 (for A,
$0.30 \times 8 + 0.10 \times 3 + 0.05 \times 1 = 2.75$). Slope $-1.194$ points per thousand
dalens of exposure. % output: ex12_shift_share
It is causal if the industry shocks are as good as randomly assigned with respect
to other regional shocks (and there are many of them), or if the initial shares
are unrelated to other causes of the regions' later employment --- for example,
if regions specialised in textiles were not already declining for other reasons.
\end{solution}

\begin{solution}{pb:b3:trade-policy:steel}
\textbf{1.} $M(P) = (100 - 2P) - (-20 + 2P) = 120 - 4P$. % figdata: steel_cases, grain_cases, payoffs

\textbf{2.} $120 - 4P = -40 + 4P$: $P = 20$, $M = 40$ (production 20, consumption 60).

\textbf{3.} $P = 25$: production 30, consumption 50, imports 20.

\textbf{4.} Consumers lose $5 \times (60 + 50)/2 = 275$; producers gain $5 \times (20 + 30)/2 = 125$;
revenue $5 \times 20 = 100$.

\textbf{5.} Deadweight loss $\tfrac12 \times 5 \times 10 + \tfrac12 \times 5 \times 10 = 50$; net $-275 + 125 + 100
= -50$.

\textbf{6.} The world price does not move, so there is no terms-of-trade gain to
set against the efficiency loss.

\textbf{7.} $120 - 5P^* = -40 + 4P^*$: $P^* = 160/9 = 17.78$, $P = 22.22$, $M = 31.11$.

\textbf{8.} The wedge is $22.22 - 17.78 = 4.44$; Kestrian exporters' price fell by 2.22,
half of it. Brenland's buyers pay the other half.

\textbf{9.} Terms-of-trade gain $(20 - 17.78) \times 31.11 = 69.1$; efficiency loss $\tfrac12 \times
(22.22 - 20) \times (40 - 31.11) = 9.9$.

\textbf{10.} $69.1 - 9.9 = 59.3$. Check: consumers $-128.4$, producers $+49.4$, revenue $4.44
\times 31.11 = 138.3$; sum 59.3.

\textbf{11.} Kestrian producer surplus falls by $(20 - 17.78) \times (40 + 31.11)/2 = 79.0$.

\textbf{12.} $\varepsilon^* = 4P^*/X^*$: $4 \times 20/40 = 2$ at free trade, $4 \times 17.78/31.11 = 2.29$ at the
tariff. $1/2.29 = 0.44 > 0.25$: the 25~\% tariff is below the optimum.

\textbf{13.} $t = 1/\varepsilon^*$ with $P^* = 40/(2 + t)$ from market clearing: $t = (4P^* - 40)/(4P^*) =
1 - (2 + t)/4$, so $t = 0.4$. Then $P^* = 16.67$, $P = 23.33$, $M = 26.67$, and Brenland gains
66.7.

\textbf{14.} $-20 + 2p = 100 - 3p$: $p = 24$, exports 28.

\textbf{15.} $-20 + 2p = 100 - 3.75p$: $p = 20.87$, $p_K = 26.09$, exports 21.74.

\textbf{16.} Brenland's grain producers lose $(24 - 20.87) \times (28 + 21.74)/2 = 77.9$.

\textbf{17.} Terms-of-trade gain $(24 - 20.87) \times 21.74 = 68.1$, minus deadweight loss
$\tfrac12 \times (26.09 - 24) \times (28 - 21.74) = 6.5$: $+61.5$.

\textbf{18.} With Brenland's payoff first: (free, free) $0, 0$; (free, tariff) $-77.9,
61.5$; (tariff, free) $59.3, -79.0$; (tariff, tariff) $-18.6, -17.5$.

\textbf{19.} A tariff is each country's dominant strategy ($59.3 > 0$ and $-18.6 > -77.9$
for Brenland; $61.5 > 0$ and $-17.5 > -79.0$ for Kestria): both levy tariffs, and both
are worse off than under free trade, a prisoner's dilemma.

\textbf{20.} In 2018--2019 the world prices of targeted U.S. imports did not fall:
foreign export supply behaved as if perfectly elastic, and the terms-of-trade
gain of Part II did not appear. % ledger: b3tp:arw.passthrough, b3tp:fgkk.loss

\textbf{21.} Steel's price rises by $22.22/20 - 1 = 11.1~\%$. ERP $= (0 - 0.3 \times 0.111)/0.7 =
-4.8~\%$.

\textbf{22.} Steel output rises from 20 to 24.44 million tonnes: about 4{,}440 jobs.

\textbf{23.} $50{,}000 \times 4.8~\% \approx 2{,}380$ jobs.

\textbf{24.} Jobs moved between industries are not a gain or a loss of welfare in
themselves; welfare depends on the surpluses, revenue and \omterm{def:b3:ricardian-trade:terms-of-trade}{terms of trade} computed
above, and workers who move find other jobs (with real costs of moving that this
count omits).

\textbf{25.} For: steel producers and their workers, whose gain is concentrated.
Against: car makers and other steel users, and grain farmers exposed to
retaliation; consumers lose most in total but are dispersed.

\textbf{26.} The tariff raises Brenland's welfare by 59.3 million dalens without
retaliation and lowers it by 18.6 million dalens with Kestria's retaliation.
\end{solution}
